\subsection{Introduction}

In the previous chapters in this thesis, we have seen a collection of different evaluations that probe different capabilities
of language models and reveal various insights about their behavior. Shaped by the lessons that we have learned from these
evaluations, we now briefly situate the task space of AI for software engineering,
then use formal verification as a representative task for discussing challenges and
paths forward.

Our goal is threefold. In Sec. \ref{sec:tasks-milestones}, we briefly summarize the
main task families in AI for software engineering and focus on formal verification as
a checkable target for reasoning about software correctness. The original position
paper contains the broader taxonomy of software-engineering tasks beyond code
generation and code completion, including the measures used to categorize their scope,
logical complexity, and level of human intervention.

Moving forward to Sec. \ref{sec:challenges}, we highlight several challenges in the field that today's models face,
each cross-cutting and applicable to several tasks. These challenges synthesize the lessons from the earlier chapters:
future AI software engineers should be evaluated across broader tasks, with metrics that preserve real
development difficulty, and in settings that account for feedback, code quality, and tool use.

In Sec. \ref{sec:paths}, we posit a set of promising research directions to tackle the challenges above.
We hope that through this chapter, the reader can appreciate the progress the field has made, understand
the shortcomings of today's state-of-the-art models, and take inspiration from our suggested future ideas
for tackling these challenges.
